% 2007 年重庆卷理科第 22 题：椭圆 x^2/36+y^2/27=1、右焦点 F(3,0)、右准线 l:x=12，
%   P1、P2、P3 为椭圆上三点且两两与 F 的连线夹角都是 120°。
% 来源 PDF：gaokaomath-src/普通高考/2007/2007重庆理.pdf
% SHA-256：83bfde9cfb8071229219565245c08cf4481943489659cb0f050919c380c65c3e
% 原页：2007重庆理.pdf 第 2 页（A3 三栏）第 3 栏，书页码 356。
% 重绘依据：以焦点为极点取 r=4.5/(1+0.5cosθ)，θ 取 30°、150°、270° 得三点
%   P1(5.72,1.57)、P2(-3.88,3.97)、P3(3,-4.5)，与原页三点方位一致。
\begin{tikzpicture}[x=0.45cm, y=0.45cm, semithick,
  every node/.style={inner sep=1.5pt}]

  % 坐标轴
  \draw[->,>=stealth] (-8.4,0) -- (14.4,0) node[right] {$x$};
  \draw[->,>=stealth] (0,-7.4) -- (0,7.4) node[above] {$y$};
  \node[below left] at (0,0) {$O$};

  % 椭圆 x^2/36+y^2/27=1
  \draw[line width=1pt] (0,0) ellipse (6 and 5.196);

  % 右准线 l：x=12
  \draw (12,-6.6) -- (12,6.6) node[above] {$l$};

  \coordinate (F)  at (3,0);
  \coordinate (P1) at (5.72,1.57);
  \coordinate (P2) at (-3.876,3.97);
  \coordinate (P3) at (3,-4.5);

  \draw (F) -- (P1);
  \draw (F) -- (P2);
  \draw (F) -- (P3);

  \fill (F) circle (1.4pt);
  \node[below right=1pt] at (F) {$F$};
  \node[above right] at (P1) {$P_1$};
  \node[above left] at (P2) {$P_2$};
  \node[above left=2pt] at (P3) {$P_3$};
\end{tikzpicture}
